\subsection{Named Entity Recognition}
\todo{add evaluation metrics. also for nel,...}
\label{sec:rw:ner}
\begin{definition}
    [Entity Recognition (\acs{ner})]
    \Acf{ner} is the task of identifying and classifying named entities or concepts in text into predefined categories such as persons, organizations, and locations~\cite{entityorientedsearch2018balog,keraghel2024recentadvancesnamedentity,jm3}.
\end{definition}

\begin{figure}[htbp]
\centering
\begin{minipage}[t]{0.48\linewidth}
\small
\PERSON{Giovanni Bianchi}, residing in \LOCATION{Bologna} \texttt{[...]} -- plaintiff\\
and\\
\PERSON{Davide Ricci}, residing in \LOCATION{Modena} \texttt{[...]} -- defendant\\
FACTS\\
The plaintiff summoned the defendant to court, alleging that on \DATE{April 3, 2022}, he lent him the sum of \MONEY{EUR 10,000}, as evidenced by a private written agreement signed by both parties, with the obligation to repay by \DATE{December 31, 2022}, pursuant to \LAW{Article 1813 c.c.} \texttt{[...]}
\subcaption{\ac{ner} on a judgment}
\label{fig:extract-ner-judgment}
\end{minipage}
\hfill
\begin{minipage}[t]{0.48\linewidth}
\small
\DATE{[2025/04/22 -- 11:45]} \PERSON{Giovanni Bianchi}: Meet me \DATE{Friday, 19:00}, at \LOCATION{Bar Centrale}. Bring the papers.\\
\DATE{[2025/04/22 -- 11:46]} \PERSON{Davide Ricci}: Fine.\\
\DATE{[2025/04/22 -- 18:51]} \PERSON{Giovanni Bianchi}: \textit{[Voice message -- 2 min 58 sec]}
\quad \textit{Transcription:}\\
\quad \texttt{[...]} You spent my money on that weekend in \LOCATION{Milan}, dinner at \LOCATION{Grand Hotel et de Milan} and even those tickets for the race at the \LOCATION{Monza Circuit} \texttt{[...]}\\
\subcaption{\ac{ner} on chat logs}
\label{fig:extract-ner-chat}
\end{minipage}
\caption{\Ac{ner} application for \acl{sta} with entity highlights and labels.}
\label{fig:extracts-ner}
\end{figure}

An example application of \ac{ner} on legal judgments and chat logs is shown in Figure~\ref{fig:extracts-ner}, where entities such as persons, locations, dates, monetary values, and legal references are identified and classified.


While traditionally focused on entities like persons, organizations, and locations, \ac{ner} systems have been applied for recognizing wide range of entity types, such as dates, monetary values~\cite{muc6_overview_1995}, biomedical entities~\cite{hanisch2005prominer,archana2023effective,aioner_bioner_2023}, legal references~\cite{chalkidis_extracting_contracts_2017,cardellino_etal_2017_lowcost,ccetindaug2023name_legal_turkish}, and more~\cite{entityorientedsearch2018balog}. The \ac{nlp} community manifested interest in extending the types of entities recognized, proposing hundreds of entity types~\cite{sekine2002extended_150ner_types,sekine2004definition_200ner_types,yosef-etal-2012-hyena} and introducing the \emph{fine-grained entity typing} field, which aims to classify entities into a larger set of more specific categories~\cite{fget_review_2023}.

\todo{maybe remove the repetition useful for sta and kbp as ie}
\ac{ner}, as an \acl{ie} task, contributes to \acl{sta}, \acl{kbp} and downstream tasks such as entity-based search or maintaining \acp{kb} up-to-date~\cite{entityorientedsearch2018balog}. For example, in \ac{sta}, recognizing entities and their allows to semantically highlight documents to let readers quickly grasp the main topics discussed from the entities mentioned. 

Besides, \ac{ner} is applied in specific domains, such as the healthcare~\cite{hanisch2005prominer,archana2023effective,aioner_bioner_2023} or the legal~\cite{chalkidis_extracting_contracts_2017,cardellino_etal_2017_lowcost,ccetindaug2023name_legal_turkish} ones, to extract domain-specific entities and support specialized applications.
Beyond these areas, \ac{ner} can contribute to document summarization~\cite{khademi2020persian,lin-zeldes-2024-gumsley,roha_ner_summarizaton_2023} and machine translation~\cite{xie2022end,hkiri-etal-2017-arabic}, where it helps preserve and highlight key entities, or question answering~\cite{molla2006named,ner_qa_review_2016}, by identifying entities mentioned in the query.

Formally, given a sequence of tokens $X$,
\begin{equation}
    X = [x_0, x_1, \ldots, x_{|X|}]
\end{equation}
and a set of categories $T$,
\begin{equation}
T = \{\texttt{Person}, \texttt{Location}, \texttt{Organization}, \ldots\},
\end{equation}
a \ac{ner} system produces a set of tuples $Y$:
\begin{equation}
Y = \{(i_{start},i_{end},t^i),\ldots\}, \text{ where } i_{start}, i_{end} \in \{0,1,\ldots,|X|\}, t^i \in T.
\end{equation}
Considering the example sentence ``Barack Obama was born on Aug 04, 1961.'', a \ac{ner} system should identify two mentions of entities: ``Barack Obama'' as a \texttt{Person} and ``Aug 04, 1961'' as a \texttt{Date}. Supposing the tokenization produces the sequence $X^I$,
\begin{equation}
X^I = [\text{``Barack''}_0, \text{``Obama''}_1, \text{``was''}_2, \text{``born''}_3, 
       \text{``on''}_4, \text{``Aug''}_5, \text{``04''}_6, \text{``,''}_7, 
       \text{``1961''}_8, \text{``.''}_9]
\end{equation}
the expected output is as follows:
\begin{equation}
Y^I = \{(0, 1, \texttt{Person}), (5, 8, \texttt{Date})\}.
\end{equation}
The first tuple correspond to ``Barack Obama'' (\texttt{Person})---from token 0 to token 1, while the second to ``Aug 04, 1961'' (\texttt{Date})---from token 5 to token 8.

Typically, \ac{ner} systems are modeled as sequence labeling tasks~\cite{jm3,keraghel2024recentadvancesnamedentity,entityorientedsearch2018balog}, where---in the simplest case---each token $x_i$ in the input sequence $X$ is assigned a label $l_i \in L$, where $L$ is defined as follow:
\begin{equation}
    T = \{\texttt{I-PER}, \texttt{I-LOC}, \texttt{I-ORG}, \texttt{O}\},
\end{equation}
\texttt{I-}\textit{TYPE} indicates that the token is \texttt{I}nside an entity of type \texttt{TYPE}, while \texttt{O} indicates that the token is \texttt{O}utside any entity mention. The output of the \ac{ner} system is thus a sequence of labels $L$, that for the above example $X^I$ would be:
\begin{equation}
L^I = [\texttt{I-PER}, \texttt{I-PER}, \texttt{O}, \texttt{O}, \texttt{O}, \texttt{I-DATE}, \texttt{I-DATE}, \texttt{O}, \texttt{I-DATE}, \texttt{O}].
\end{equation}
The \emph{IO} scheme is the simplest labeling scheme for \ac{ner} and it cannot distinguish between consecutive entities of the same type~\cite{alshammari2021impact}. For example, consider the sentence ``Barack Obama Michelle Obama'' from the keywords of a newspaper, the IO scheme would label it as:
\begin{equation}
L^{II} = [\texttt{I-PER}, \texttt{I-PER}, \texttt{I-PER}, \texttt{I-PER}],
\end{equation}
thus failing to recognize the two distinct mentions of \texttt{Person}. More sophisticated and informative schemes exists, such as the \emph{BIO}~\cite{ramshaw-marcus-1995-text-bio-scheme}, which introduces the \texttt{B-}\textit{TYPE} label to indicate that the token is at the \texttt{B}eginning of an entity of type \texttt{TYPE}. With BIO, the above example would be labeled as:
\begin{equation}
L^{II} = [\texttt{B-PER}, \texttt{I-PER}, \texttt{B-PER}, \texttt{I-PER}],
\end{equation}
thus correctly distinguishing two mentions of \texttt{Person}.
Refer to \textcite{alshammari2021impact} for a comparison of different annotation schemes for \ac{ner} and their impact on model performance.

Early \ac{ner} systems were based on rule-based or manually engineered features, rules, or gazetteers. These approaches typically achieved high precision but low recall~\cite{keraghel2024recentadvancesnamedentity}.
Later, machine learning methods enabled more adaptable and data-driven approaches. Effective \ac{ner} models have been built using hidden Markov models, support vector machines, and conditional random fields (CRFs)~\cite{keraghel2024recentadvancesnamedentity}---able to consider entire sentences for \ac{ner} predictions.
With deep learning, and representation learning, \ac{ner} systems further improved. Combinations of deep architectures with CRFs---e.g., BiLSTM-CRF~\cite{riedl-pado-2018-named,schweter-baiter-2019-towards} or BiLSTM-CNN-CRF~\cite{ma-hovy-2016-end}---achieved strong results~\cite{keraghel2024recentadvancesnamedentity}. The introduction of transformers~\cite{vaswani_transformers_2017}, BERT~\cite{devlin-etal-2019-bert} and its variants, reshaped the field, outperforming previous BiLSTM-CRF models on standard benchmarks~\cite{keraghel2024recentadvancesnamedentity}.




More recently, \ac{ner} has been approached using \aclp{llm} with different techniques, such as zero-shot prompting, \acl{icl}, and supervised fine-tuning~\cite{xu2024llm_for_ie}. For example, \textcite{wang-etal-2025-gptner} performs \ac{ner} via \ac{icl} by mapping the task to a text generation problem~\cite{}---i.e., by asking the model to repeat the input text surrounding entity mentions with a prefix and suffix (e.g., ``<PERSON> Barack Obama </PERSON>''). Other approaches prompt \acp{llm} to enumerate entity mentions in a structured format without repeating the input text~\cite{xie-etal-2024-self}, or even to produce code that represents the extracted entities~\cite{sainz2024gollie,li-etal-2023-codeie}---for instance, by defining a class for each entity type with fields such as the mention string, and then asking the model to output a list of entity objects~\cite{sainz2024gollie}.

The \ac{llm}-based approaches have shown strong adaptability, thanks to their zero-shot and few-shot capabilities, reducing the need for annotating large datasets for domain-specific applications~\cite{keraghel2024recentadvancesnamedentity}. However, \acp{llm} are resource intensive and their adoption is limited by infrastructural costs and latency~\cite{keraghel2024recentadvancesnamedentity}. Also designing effective prompts is an expensive process---requiring domain expertise---and small changes in the prompt can vary the model's behavior, making these approaches not suitable for domains where consistent outputs are required~\cite{keraghel2024recentadvancesnamedentity}. 

Recent studies have explored the use of smaller models---i.e., DeBERTa~\cite{HeGChen2023} and RoBERTa~\cite{liu2020roberta}, which are improvement of BERT~\cite{devlin-etal-2019-bert}---for zero-shot \ac{ner}, where the entity types to recognize are provided as textual inputs for the model together with the input text~\cite{zaratiana-etal-2024-gliner,bogdanov-etal-2024-nuner}. An example application is to recognize entities of type \texttt{Clothing} or \texttt{Pet} by proving the strings ``clothing'' and ``pets'' in the input. These models achieve competitive performance compared to \acp{llm}, while being more efficient and easier to deploy~\cite{zaratiana-etal-2024-gliner,bogdanov-etal-2024-nuner}.

\todo{add uniner}

For further information on recent advancements, refer to \textcite{keraghel2024recentadvancesnamedentity} for \ac{ner} and \textcite{xu2024llm_for_ie} for \ac{ie}---including \ac{ner}---with \acp{llm}.

\todo{order: definition, example figure, literature and advancements, "production" libraries eg spacy, tint,...}
\todo{ner for italian (and legal)?}

Finally, several ready-to-use libraries exist for \ac{ner}, such as \spacy~\cite{honnibal2020spacy}, Flair~\cite{akbik2019flair}, Stanza~\cite{qi2020stanza}, Stanford CoreNLP~\cite{manning-EtAl:2014:P14-5}, and Transformers~\cite{wolf-etal-2020-transformers}. These libraries provide pre-trained models for multiple languages, and can be easily integrated into applications. For instance, \spacy offers efficient and accurate \ac{ner} models that can be fine-tuned on custom datasets~\cite{honnibal2020spacy}. Its models are based on a transition-based parser~\cite{AliEtAl2023UnintendedMemorization} that maps structured prediction tasks, such as \ac{ner} and dependency parsing, into a sequence of state transitions~\cite{lample-etal-2016-neural}. The underlying architecture combines multiple subnetworks. The tok2vec component encodes tokens into embeddings enriched with sub-word features (e.g., prefixes, suffixes, or shapes such as ``xxxXx'' for ``spaCy'') using a CNN-based encoder. A lower subnetwork then constructs representations for each token by combining these features, while an optional upper feed-forward network refines the representation to predict action scores. When the upper network is not used, the lower subnetwork directly produces the scores~\cite{AliEtAl2023UnintendedMemorization}. SpaCy also provide transformer-based models in which the tok2vec component is replaced by a transformer---for example the ``en\_core\_web\_trf'' model\footnote{\url{https://spacy.io/models/en\#en\_core\_web\_trf}}.

\todo{how to perform ner with BERT, training objective, eval metrics}




\todo{nested entities, overlapping entities}